\documentclass{article}

\usepackage{../project-format/project_440_550}
\usepackage[utf8]{inputenc}
\usepackage[T1]{fontenc}
\usepackage[USenglish]{babel}
\usepackage{csquotes}
\usepackage{booktabs}
\usepackage{amsfonts}
\usepackage{amsmath}
\usepackage{nicefrac}
\usepackage{microtype}
\usepackage{xcolor}
\usepackage{hyperref}
\usepackage{url}
\usepackage[round]{natbib}
\usepackage[capitalize,noabbrev]{cleveref}

\title{Eigenslang: A Geometric Approach to Youth Slang Normalization}

\author{%
  Yuhei Arimoto
  \And
  James Wen
}

\begin{document}
\maketitle

\begin{abstract}
As LLMs become increasingly integrated into the lives of young users, the rapid evolution of Gen Z and Gen Alpha slang has created a critical misalignment between NLP systems and their primary demographic. Current normalization models often rely on opaque ``black-box'' architectures that lack interpretability. We hypothesize that the semantic shift between youth slang and standard English is not arbitrary, but is instead partly defined by identifiable, low-dimensional directions in embedding space, which we term \emph{Eigenslang}. We evaluate this hypothesis using a curated dataset of 96 slang--neutral pairs, comparing mean-offset and PCA-based methods against raw retrieval baselines. Our results provide only weak evidence for a single shared slang direction, with the first principal component explaining only about 5\% of the variance.
An abbreviation-filtered ablation improves raw retrieval substantially, which suggests that dataset heterogeneity matters, but it still does not produce strong evidence for a single global PCA direction. We conclude that while representation choice significantly affects performance, a single global direction is an incomplete explanation of the complex nature of slang normalization.
\end{abstract}

\section{Introduction}
Modern NLP systems still struggle with youth slang, especially with the rapidly evolving Gen Z and Gen Alpha language prevalent on social media. Common Gen Z expressions such as \textit{mid}, \textit{cooked}, \textit{ate}, and \textit{rizz} are recent, informal, and context-dependent, so standard lexical resources and embeddings often represent them poorly. Despite strategies such as subword tokenization, contextual embeddings, and more diverse training data, challenges remain due to the fluid and rapidly evolving nature of slang \citep{ishita-mamidi-2025-evolution}. Furthermore, existing normalization frameworks, such as neural machine translation systems, typically rely on ``black-box'' encoder-decoder architectures. These models lack interpretability, providing no clear insight into the semantic transformations they apply to informal text. This is detrimental for tasks such as text understanding, retrieval, sentiment analysis, and chatbot interaction. Ultimately, the inability of these systems to accurately and transparently parse slang creates a fundamental disconnect in user experience and accessibility.

In recent years, there have been significant advances in computational linguistics. Specifically, the application of eigenvector analysis has emerged as a powerful tool for uncovering latent semantic structures and providing a more robust interpretation of high-dimensional word embeddings. This leads us to the main research question of our project: does the \emph{Eigenslang} exist? That is, is there a structured transformation from a neutral phrase to its slang counterpart? We call the shared direction, if it exists, the \emph{Eigenslang}. If a slang term and a neutral paraphrase are embedded as vectors, is the difference between them partly shared across many examples? If so, that direction could represent ``slanginess'' and support a lightweight translation or normalization method.

Prior work such as \textit{Interpreting Word Embeddings with Eigenvector Analysis} by \citet{shin2018eigenvectors} suggests that eigenvector-based analysis can reveal interpretable structure in embedding spaces. Other work on slang extraction, slang generation, Gen Alpha translation, and youth-slang chatbots further motivates this as both a linguistic and ML problem \citep{ishita-mamidi-2025-evolution}. We also drew inspiration from the theoretical concept and framework of the \textit{"Eigenslur"}, popularized by linguistic influencers like 
etymologynerd.

The scope of this project is intentionally focused and modest. Rather than claiming a universal solution to slang normalization, we provide a controlled empirical evaluation of whether a low-dimensional structure exists within slang-neutral offsets. By isolating these semantic directions, we aim to offer a more interpretable alternative to traditional 'black-box' machine translation methods.

% This goes in conclusion
% Regardless of the results, both positive and negative results are meaningful:
% if the leading direction helps retrieval, that supports the hypothesis;
% if not, that suggests slang may be too heterogeneous or context-specific for a single global offset.

% Our contribution is a curated 96-example dataset, a comparison of two offset formulations, and a controlled test showing that representation choice and dataset composition matter more than a single shared PCA direction.

\section{Related Work}
Our work is motivated by three strands of literature.

First, prior work shows that slang benefits from specialized lexical resources and embeddings.
\citet{wilson2020urban} build Urban Dictionary-based embeddings for slang NLP and show that slang-specific lexical information can improve downstream behavior.
Similarly, social-media slang extraction work argues that non-standard language behaves differently from standard lexicons and often requires specialized treatment \citep{matsumoto2019slang}.

Second, our project is motivated by recent advancements in identifying interpretable structure within embedding spaces. \citet{shin2018eigenvectors} show that eigenvector-based analysis of word embeddings can uncover latent semantic organization that is otherwise difficult to inspect directly. More recently, \citet{shaik2026redefining} use eigen-decomposition to study and intervene on harmful behavior in language models, illustrating a broader interpretability trend toward identifying structured components rather than relying only on opaque end-to-end behavior.

This aligns with the broader consensus in the field that meaning can be organized along approximately linear directions \citep{bojanowski2017enriching, bolukbasi2016man}.
Our project adopts this geometric style of reasoning, extending it to the domain of slang normalization. Unlike previous work focused on removing demographic bias or solving classic analogies, we investigate whether this linear structure can be harnessed to map the shift from youth vernacular to standard English in a transparent, interpretable manner.

Third, there is prior work on slang generation and transformation.
\citet{sun2021computational} propose a computational framework for slang generation.
Our problem is simpler and more controlled: we study retrieval-based normalization rather than generation.

FastText is a natural point of comparison because subword structure can help with rare or novel forms \citep{bojanowski2017enriching}.
In the current implementation, we instead compare TF--IDF and a sentence-transformer model based on Sentence-BERT \citep{reimers2019sentence}, because many of our examples are phrase-level or context-dependent.
This design choice became especially important once we observed that isolated word-level paraphrases and full context sentences can behave differently.

\section{Dataset and Task}
\subsection{Dataset Construction}
We began with a small manually written dataset and then expanded it using a cleaned subset of the publicly available \texttt{MLBtrio/genz-slang-dataset} on Hugging Face \citep{mlbtrio2026dataset}.
The external dataset was used as a source pool rather than as trusted final data.
We imported raw rows, mapped them into our schema, filtered duplicates and low-quality entries, and selected external candidates conservatively.

The resulting cleaned dataset contains 96 examples and serves as our primary evaluation set: 60 manual examples and 36 selected external examples, with fields for slang term, context sentence, neutral expression, neutral paraphrase, and neutral sentence.
It also stores coarse category labels such as \texttt{judgment}, \texttt{emotion}, \texttt{online}, and \texttt{dating\_social}.

\subsection{Why We Added Neutral Fields}
One neutral target is not enough for this task, because the same example is used in two different ways.
For direct term/paraphrase experiments, we want a short canonical normalization target.
For contextual experiments, we want a full sentence that reads naturally after the slang expression is removed.
These are not always the same object.

For example, a descriptive gloss such as \emph{received far more criticism than support} is a reasonable semantic target, but inserting it directly into a sentence can produce awkward text.
We therefore store three neutral-side fields:
\begin{itemize}
\item \texttt{neutral\_expression}: a short neutral word or phrase used for direct term/paraphrase offsets;
\item \texttt{neutral\_paraphrase}: a fuller descriptive gloss that captures the intended meaning;
\item \texttt{neutral\_sentence}: a manually neutralized sentence used for context-preserving comparisons.
\end{itemize}

This distinction lets us evaluate both lexical and contextual formulations without forcing one neutral representation to serve both purposes poorly.

\subsection{Task Definition}
We formulate slang normalization as retrieval.
Given a slang-side representation, the system must retrieve the correct neutral target from a candidate set.
We report Top-1 and Top-3 retrieval accuracy.

All three evaluation settings use the same underlying dataset.
What changes across settings is not the dataset itself, but the set of candidate neutral targets presented for each test example.

We report three candidate-set settings.
In \texttt{all\_candidates}, the correct neutral target is ranked against all neutral labels in the dataset.
In \texttt{category\_candidates}, the candidate pool is restricted to the example's category.
In \texttt{balanced\_distractors}, we construct a smaller controlled candidate set consisting of the correct neutral target plus a fixed-size sample of distractors, rather than the full dataset.
This creates a setting that is still competitive, but less dominated by brute-force search over all labels.
These settings let us distinguish between semantic failure and evaluation difficulty.

\section{Methods}
\subsection{Two Offset Formulations}
We define the basic difference vector as
\[
\mathbf{d}_i = \mathbf{s}_i - \mathbf{n}_i,
\]
where $\mathbf{s}_i$ is a slang embedding and $\mathbf{n}_i$ is the embedding of its neutral paraphrase.

In implementation, we evaluate two versions of this idea.

\paragraph{Term/paraphrase offset.}
This is the direct formulation:
\[
\mathbf{d}_i =
\operatorname{Embed}(\text{slang term})
-
\operatorname{Embed}(\text{neutral expression}).
\]
This formulation is simple and directly tests whether slang and neutral expressions lie along a common offset in embedding space.

\paragraph{Contextual sentence offset.}
Because many slang items are polysemous or phrase-level, we also evaluate
\[
\mathbf{d}_i =
\operatorname{Embed}(\text{slang sentence})
-
\operatorname{Embed}(\text{neutral sentence}).
\]
For example, the standard and slang meanings of \emph{ate}, \emph{cooked}, and \emph{serve} can differ substantially.
This contextual formulation attempts to preserve the intended meaning in a way that isolated words may not.

\subsection{Embedding Techniques}
We compare two embedding techniques:
\begin{itemize}
\item TF--IDF, as a simple lexical baseline;
\item a sentence-transformer model using \texttt{all-MiniLM-L6-v2} \citep{reimers2019sentence}.
\end{itemize}

TF--IDF is not expected to capture slang semantics well, but it is useful as a sanity check.
The sentence-transformer is a stronger semantic baseline and supports both phrase-level and sentence-level inputs.

\subsection{Baselines and Eigenslang}
For each training example we compute a slang--neutral difference vector.
From these differences we derive three methods:

\paragraph{Raw retrieval.}
We compare the slang-side representation directly to all neutral candidates using cosine similarity.

\paragraph{Mean offset.}
We compute the mean training difference vector and assume it as the candidate Eigenslang.
\[
\bar{\mathbf{d}} = \frac{1}{N} \sum_i \mathbf{d}_i
\]
and apply
\[
\mathbf{s}' = \mathbf{s} - \alpha \bar{\mathbf{d}}.
\]

\paragraph{PCA Eigenslang.}
We perform PCA on the training difference vectors and take the first principal component $\mathbf{e}_1$ as the candidate Eigenslang.
We then apply
\[
\mathbf{s}' = \mathbf{s} - \alpha \mathbf{e}_1.
\]
The scaling factor $\alpha$ is tuned on validation data.

\section{Experiments}
\subsection{Setup}
We split the cleaned dataset into train, validation, and test sets with a deterministic 60/20/20 split. We evaluate raw retrieval, mean offset, and PCA Eigenslang under both representation modes and both embedding techniques. At test time, the model starts from the slang-side representation and retrieves the correct neutral target from a candidate set.

\textbf{Top-1 Accuracy}: The percentage of cases where the correct neutral target is ranked first.

\textbf{Top-3 Accuracy}: The percentage of cases where the correct neutral target appears in the top three predictions.

This setup measures how well a learned offset helps move slang representations toward their intended neutral meanings.

\subsection{Main Quantitative Results}
\cref{tab:candidate-results} summarizes the sentence-transformer test-set results on the primary 96-example dataset across candidate-set settings.
We focus on sentence-transformers in the main table because TF--IDF is uniformly near zero and does not affect the qualitative conclusions.

\begin{table}
  \caption{Sentence-transformer test accuracy under different candidate sets.}
  \label{tab:candidate-results}
  \centering
  \begin{tabular}{llllll}
    \toprule
    Representation & Candidate set & Method & Top-1 & Top-3 & $\alpha$ \\
    \midrule
    term/paraphrase & all & raw & 0.10 & 0.25 & -- \\
    term/paraphrase & all & mean offset & 0.10 & 0.25 & 0.25 \\
    term/paraphrase & all & PCA Eigenslang & 0.10 & 0.25 & 0.00 \\
    term/paraphrase & category & raw & 0.35 & 0.65 & -- \\
    term/paraphrase & category & mean offset & 0.35 & 0.70 & 1.50 \\
    term/paraphrase & category & PCA Eigenslang & 0.35 & 0.65 & 0.00 \\
    term/paraphrase & balanced & raw & 0.40 & 0.80 & -- \\
    term/paraphrase & balanced & mean offset & 0.45 & 0.75 & 0.50 \\
    term/paraphrase & balanced & PCA Eigenslang & 0.40 & 0.80 & 0.00 \\
    contextual & all & raw & 0.00 & 0.05 & -- \\
    contextual & all & mean offset & 0.00 & 0.05 & 0.00 \\
    contextual & all & PCA Eigenslang & 0.05 & 0.10 & 0.25 \\
    contextual & category & raw & 0.30 & 0.50 & -- \\
    contextual & category & mean offset & 0.20 & 0.55 & 2.00 \\
    contextual & category & PCA Eigenslang & 0.15 & 0.50 & 0.25 \\
    contextual & balanced & raw & 0.40 & 0.70 & -- \\
    contextual & balanced & mean offset & 0.30 & 0.60 & 2.00 \\
    contextual & balanced & PCA Eigenslang & 0.30 & 0.65 & 0.25 \\
    \bottomrule
  \end{tabular}
\end{table}

\subsection{What the Results Suggest}
Firstly, we find that TF--IDF is clearly too weak for this task.
This is not surprising, since many of the examples depend on semantic nuance rather than surface-level lexical overlap.

Second, the direct term/paraphrase formulation with a sentence-transformer model gives the best absolute performance.
Under the more controlled candidate sets, it reaches Top-1 = 0.35 and Top-3 = 0.70 in the category-restricted setting, and Top-1 = 0.45 and Top-3 = 0.75 with balanced distractors.
These results are much stronger than the original all-candidate numbers, which shows that evaluation design matters substantially.
However, PCA Eigenslang still does not improve over raw retrieval in this formulation.

Third, the contextual sentence formulation shows the only small positive effect from PCA Eigenslang, but only in the harshest all-candidate setting.
With sentence-transformers, raw retrieval achieves Top-1 = 0.00 and Top-3 = 0.05 in that setting, while PCA reaches Top-1 = 0.05 and Top-3 = 0.10.
Once we move to category-restricted or balanced candidate sets, that PCA advantage disappears and raw retrieval is again stronger.

\subsection{PCA Variance Analysis}
To test the Eigenslang hypothesis more directly, we also measured how much variance the first principal component explains in the training difference vectors.
For the sentence-transformer runs with category-restricted candidates, the first component explains 5.28\% of the variance in the term/paraphrase formulation and 5.45\% in the contextual formulation.
The top three components explain only 13.77\% and 15.24\%, respectively.

These numbers are important because they show that the slang--neutral differences do not collapse onto one strong global axis.
If a dominant Eigenslang direction existed, we would expect the first principal component to explain much more of the total variance.
Instead, the geometry appears diffuse and multi-factor.

\subsection{Ablation on Abbreviations}
We also tested whether abbreviation-heavy items were making the dataset too heterogeneous for a single shared slang direction.
This ablation was run on a separate merged dataset created by augmenting the primary 96-example dataset with 16 curated examples from the \emph{Gen Z words and Phrases Dataset} \citep{yeasmin2026kaggle}, yielding 112 rows in total.
From that merged dataset, we removed 27 obvious shorthand or acronym-like items, such as \emph{IRL}, \emph{SMH}, \emph{G2G}, \emph{LMAO}, \emph{BFF}, and \emph{DM}, producing an 85-row abbreviation-filtered variant.
The Kaggle-derived merged dataset did not improve the main sentence-transformer results and is therefore treated only as a robustness check.

This ablation improved raw retrieval substantially, especially in the contextual formulation, which suggests that acronym expansion and slang normalization are not quite the same phenomenon.
However, the PCA story did not change materially: even in this cleaner setting, the first principal component still explained only about 5.20\% of the variance in the term/paraphrase formulation and 6.16\% in the contextual formulation.
Detailed qualitative examples and ablation numbers are given in the appendix.

\section{Discussion and Future Work}\label{sec:discussion}
Our current conclusion is cautious.
We do \emph{not} have strong evidence that a robust, universal Eigenslang direction exists.
The direct term/paraphrase formulation performs best overall, but PCA does not improve it.
The contextual sentence formulation only shows a small PCA gain in the harshest all-candidate setting.

Our contribution is a curated 96-example dataset, a comparison of two offset formulations, and a controlled test showing that representation choice and dataset composition matter more than a single shared PCA direction.

\paragraph{Strengths.}
The strongest part of the project is methodological clarity.
We started from a simple offset hypothesis, discovered a genuine implementation issue caused by polysemy and phrase-level slang, and turned that into an explicit comparison between term/paraphrase and contextual sentence offsets.
This makes the project more rigorous than if we had assumed one representation choice was obviously correct.
The later candidate-set experiments and abbreviation-filtered ablation also made the negative result more convincing by showing that weak PCA behavior persists even after making the task cleaner and easier.

\paragraph{Weaknesses.}
The main weakness is dataset scale.
Even after cleaning and expansion, the final test split is still small.
Some neutral sentences are also only approximately natural, and category labels are somewhat coarse.
Finally, we have not yet tested fastText, which could still be informative for morphologically novel slang.

\paragraph{What we would do next.}

With more time, the next three steps would be: add a fastText baseline and compare it directly against sentence-transformers under the same candidate sets; expand the clean dataset further while keeping lexical slang and acronym-like shorthand more clearly separated; add per-category error analysis and more qualitative examples of where mean offset helps or fails.

We would still consider fastText because subword information could be useful for rare or morphologically novel slang.
At the same time, our current results suggest that evaluation design and data quality were more important bottlenecks than the absence of one particular embedding technique.

\section{Conclusion}
We tested whether youth slang normalization has a shared low-dimensional structure in embedding space.
Using a cleaned 96-example dataset, we compared TF--IDF and sentence-transformer embeddings, raw retrieval, mean-offset baselines, and PCA-based Eigenslang directions under both direct term/paraphrase and contextual sentence formulations.
We also tested a Kaggle-derived augmentation and an abbreviation-filtered ablation to see whether data composition changed the conclusion.

The current evidence is mixed.
The term/paraphrase formulation performs better overall, especially under controlled candidate sets, but PCA does not help it.
The contextual formulation shows a small PCA gain only in the harshest all-candidate setting, and that gain disappears under fairer evaluation.
Removing abbreviations improves raw retrieval substantially, which shows that dataset heterogeneity matters, but even in that cleaner setting PCA remains weak and the leading component still explains only a small fraction of the variance.
Combined with the PCA variance analysis, our best current conclusion is that a single global Eigenslang direction is at most a weak and representation-dependent phenomenon.
This is still informative: it suggests that slang normalization may require category structure, richer context, or more local transformations than a single shared offset can provide.

\bibliographystyle{plainnat}
\begin{thebibliography}{99}

\bibitem[Bolukbasi et~al.(2016)Bolukbasi, Chang, Zou, Saligrama, and
Kalai]{bolukbasi2016man}
Tolga Bolukbasi, Kai-Wei Chang, James Zou, Venkatesh Saligrama, and Adam
Kalai.
\newblock 2016.
\newblock Man is to computer programmer as woman is to homemaker? debiasing
word embeddings.
\newblock In \emph{Advances in Neural Information Processing Systems}.

\bibitem[Bojanowski et~al.(2017)Bojanowski, Grave, Joulin, and
Mikolov]{bojanowski2017enriching}
Piotr Bojanowski, Edouard Grave, Armand Joulin, and Tomas Mikolov.
\newblock 2017.
\newblock Enriching word vectors with subword information.
\newblock \emph{Transactions of the Association for Computational Linguistics},
5:135--146.

\bibitem[MLBtrio(2026)]{mlbtrio2026dataset}
MLBtrio.
\newblock 2026.
\newblock \emph{MLBtrio/genz-slang-dataset}.
\newblock \url{https://huggingface.co/datasets/MLBtrio/genz-slang-dataset}.
Accessed 2026-04-23.

\bibitem[Matsumoto et~al.(2019)Matsumoto, Ren, Matsuoka, Yoshida, and
Kita]{matsumoto2019slang}
Kazuyuki Matsumoto, Fuji Ren, Masaya Matsuoka, Minoru Yoshida, and Kenji Kita.
\newblock 2019.
\newblock Slang feature extraction by analysing topic change on social media.
\newblock \emph{CAAI Transactions on Intelligence Technology}.

\bibitem[Yeasmin(2026)]{yeasmin2026kaggle}
Tawfia Yeasmin.
\newblock 2026.
\newblock \emph{Gen Z words and Phrases Dataset}.
\newblock \url{https://www.kaggle.com/datasets/tawfiayeasmin/gen-z-words-and-phrases-dataset}.
Accessed 2026-04-24.

\bibitem[Reimers and Gurevych(2019)]{reimers2019sentence}
Nils Reimers and Iryna Gurevych.
\newblock 2019.
\newblock Sentence-BERT: Sentence embeddings using siamese BERT-networks.
\newblock In \emph{Proceedings of the 2019 Conference on Empirical Methods in
Natural Language Processing}.

\bibitem[Shin et~al.(2018)Shin, Madotto, and Fung]{shin2018eigenvectors}
Jamin Shin, Andrea Madotto, and Pascale Fung.
\newblock 2018.
\newblock Interpreting word embeddings with eigenvector analysis.
\newblock In \emph{NeurIPS Workshop on Interpretability and Robustness for
Audio, Speech, and Language}.

\bibitem[Sun et~al.(2021)Sun, Zemel, and Xu]{sun2021computational}
Zhewei Sun, Richard Zemel, and Yang Xu.
\newblock 2021.
\newblock A computational framework for slang generation.
\newblock \emph{Transactions of the Association for Computational Linguistics},
9:371--386.

\bibitem[Wilson et~al.(2020)Wilson, Magdy, McGillivray, Garimella, and
Tyson]{wilson2020urban}
Steven Wilson, Walid Magdy, Barbara McGillivray, Kiran Garimella, and Gareth
Tyson.
\newblock 2020.
\newblock Urban dictionary embeddings for slang NLP applications.
\newblock In \emph{Proceedings of the Twelfth Language Resources and Evaluation
Conference}, pages 1557--1564.

\bibitem[Shaik et~al.(2025)Shaik, Mazhar, Srivastava, and Akhtar]{shaik2026redefining}
Zuhair~Hasan Shaik, Abdullah Mazhar, Aseem Srivastava, and Md~Shad Akhtar.
\newblock 2025.
\newblock Redefining experts: Interpretable decomposition of language models for toxicity mitigation.
\newblock In \emph{Thirty-ninth Conference on Neural Information Processing Systems (NeurIPS 2025), poster}.
\newblock \url{https://openreview.net/forum?id=1wmP48quNb}

\bibitem[Ishita and Mamidi(2025)]{ishita-mamidi-2025-evolution}
Ishita and Radhika Mamidi.
\newblock 2025.
\newblock The evolution of Gen Alpha slang: Linguistic patterns and {AI} translation challenges.
\newblock In \emph{Proceedings of the 63rd Annual Meeting of the Association for Computational Linguistics (Volume 4: Student Research Workshop)}, pages 678--686.
\newblock \url{https://aclanthology.org/2025.acl-srw.43/}

\end{thebibliography}

\appendix

\section{Reproducibility Notes}
All experiments in the current draft are based on local scripts in the project repository.
Code, data-processing scripts, and experiment files are available at \url{https://github.com/yuheiarimoto616/eigenslang}.
The cleaned dataset used in the main experiments is \texttt{data/processed/slang\_examples\_clean.csv}.
The main summary table in \cref{tab:candidate-results} is based on \texttt{report\_summary\_test\_metrics.csv}.
The underlying sentence-transformer metric files and PCA spectrum files are in \texttt{results/tables/}.
The most relevant files are:
\begin{itemize}
\item \path{results/tables/baseline_metrics_sentence_transformers_all_minilm_l6_v2_term_paraphrase_all_candidates.csv}
\item \path{results/tables/baseline_metrics_sentence_transformers_all_minilm_l6_v2_term_paraphrase_category_candidates.csv}
\item \path{results/tables/baseline_metrics_sentence_transformers_all_minilm_l6_v2_term_paraphrase_balanced_distractors.csv}
\item \path{results/tables/baseline_metrics_sentence_transformers_all_minilm_l6_v2_contextual_sentence_all_candidates.csv}
\item \path{results/tables/baseline_metrics_sentence_transformers_all_minilm_l6_v2_contextual_sentence_category_candidates.csv}
\item \path{results/tables/baseline_metrics_sentence_transformers_all_minilm_l6_v2_contextual_sentence_balanced_distractors.csv}
\item \path{results/tables/pca_spectrum_sentence_transformers_all_minilm_l6_v2_term_paraphrase_category_candidates.csv}
\item \path{results/tables/pca_spectrum_sentence_transformers_all_minilm_l6_v2_contextual_sentence_category_candidates.csv}
\end{itemize}

\section{Ablation and Qualitative Details}
We also tested a small Kaggle-derived augmentation based on the \emph{Gen Z words and Phrases Dataset} \citep{yeasmin2026kaggle}, but it did not improve the main experimental story and is therefore treated as a robustness check rather than a replacement dataset.

For the abbreviation-filtered ablation, the main changes were:
\begin{itemize}
\item contextual category-candidates improved from 0.22 / 0.43 to 0.41 / 0.71 in Top-1 / Top-3;
\item contextual balanced-distractors improved from 0.39 / 0.70 to 0.65 / 0.88;
\item term/paraphrase all-candidates improved from 0.04 / 0.17 to 0.18 / 0.29.
\end{itemize}

Qualitatively, raw retrieval often already captures the intended neutral meaning.
For example, \emph{fire} retrieves \emph{excellent}, \emph{rizz} retrieves \emph{charisma}, and \emph{Normalize} retrieves \emph{Make something the norm.}
By contrast, PCA usually only reshuffles nearby candidates.
For \emph{fire}, PCA moves \emph{excellent} from second to first, which is a small improvement.
But for \emph{hits different}, PCA demotes the correct target \emph{feels especially good} below more generic positive labels, and for \emph{ick}, PCA removes the correct target from the top three entirely.

\end{document}
